\section{Matriz de indicadores mínimos de sostenibilidad}
\label{sec:matriz-sostenibilidad}
\subsection{Construcción, selección y alcance de la matriz}
La matriz operacionaliza el primer objetivo específico: reúne criterios de energía, envolvente, confort, agua, materiales y operación pertinentes para un edificio público en Galápagos. La selección responde al clima cálido-húmedo, la permanencia de usuarios, la disponibilidad de agua y las exigencias de abastecimiento y mantenimiento insular. Se incluyen indicadores medibles mediante simulación y otros que requieren planos, fichas, cómputos o verificaciones posteriores. La falta de datos no elimina un indicador relevante.

La integración se realiza en cuatro pasos: identificar el criterio y su versión de origen; comprobar su pertinencia para oficinas y clima local; definir unidad, cálculo, referencia y evidencia; y aplicar la misma ficha al CB y a cada alternativa. Cuando dos marcos tratan el mismo fenómeno, se conserva un indicador y se registran las correspondencias, sin sumar puntos duplicados. La matriz es una adaptación académica mínima para este caso; su transferencia a otros edificios requiere revisar uso, clima y normativa.

Se distinguen tres referencias: el \textbf{CB de diseño}, que permite medir cambios entre escenarios; el \textbf{requisito externo}, con el alcance de su propio sistema; y el \textbf{criterio de investigación}, adoptado explícitamente cuando no se aplica un umbral externo. Una reducción frente al CB no demuestra, por sí sola, el cumplimiento de un estándar. La matriz no reemplaza la normativa local ni constituye un sistema de certificación combinado.

\clearpage
\subsection{Documentos de referencia y control de versiones}
Se adopta Minergie LATAM 2025.1 y su guía para la variante oficinas. Su certificación exige todos los requisitos obligatorios y al menos un tercio de los electivos aplicables; los indicadores seleccionados aquí no reproducen el reglamento completo \parencite{minergieLatamReglamento2025,minergieLatamGuia2025}.

Para EDGE se utilizan las guías V3 de certificación y materiales consultadas el 12 de septiembre de 2026. Se distinguen las categorías de energía operacional, agua y materiales. La guía material utiliza carbono incorporado; no se intercambian resultados de carbono y energía incorporada \parencite{ifcEdgeCertificacion2026,ifcEdgeMateriales2025}. La referencia EDGE debe construirse en su herramienta y no se presume idéntica al CB de DesignBuilder.

Para WELL se fija la edición v2 Q4 2020, características T01 y X06; se utiliza como referencia de diseño, sin afirmar adhesión a todas las actualizaciones posteriores. Para iluminación anual se fija LEED v4.1 BD+C, guía de abril de 2019, crédito EQ Daylight, opción 1. Un crédito opcional no se presenta como requisito universal de certificación \parencite{iwbiWellV22020,usgbcLeedV41Guide2019}. CEELA aporta sus 15 principios EECA como marco de diseño y operación; no se le asigna una puntuación de certificación \parencite{ceelaPrincipiosEECA2024}.

\subsection{Reglas de evaluación}
Cada fila registra evidencia, resultado y acción. Los estados son: \textbf{C}, criterio acreditado con evidencia compatible; \textbf{NC}, criterio evaluable no alcanzado; \textbf{EP}, evidencia parcial que no permite decidir; \textbf{NE}, no evaluable con la documentación disponible; \textbf{LB}, valor de línea base sin juicio de cumplimiento; y \textbf{NA}, no aplicable con justificación. Un requisito con varias condiciones solo recibe C cuando todas están acreditadas. La ausencia de datos se conserva como EP o NE y no se convierte en incumplimiento ni en cero consumo.

No se calcula un porcentaje global de sostenibilidad: los indicadores tienen alcances distintos y no se han validado ponderaciones. La cobertura documental y el cumplimiento se informan por separado. Los criterios condicionales requieren comprobar antes su aplicabilidad. Para decidir C o NC frente a una referencia térmica o energética se deben cerrar las incidencias relevantes de cargas, geometría y sistemas del modelo.

\clearpage
\subsection{Energía y envolvente}
\begin{table}[H]\centering\caption{Indicadores mínimos: energía y envolvente}
{\fontsize{9.5}{12}\selectfont\singlespacing\renewcommand{\arraystretch}{1.25}\setlength{\tabcolsep}{4pt}\begin{tabular}{>{\raggedright\arraybackslash}p{2.9cm}>{\raggedright\arraybackslash}p{5.2cm}>{\raggedright\arraybackslash}p{5.9cm}}\toprule Indicador & Criterio de referencia & Cálculo, alcance y fuente \\ \midrule
E01 · Demanda de enfriamiento & Reducir respecto al CB; criterio comparativo propio, sin umbral de certificación. & Integrar potencia horaria por zona. Ahorro = 100 × (CB $-$ caso) / CB. Mismas zonas, clima, calendario y cargas; no sumar enfriamiento sensible y total.\par\textit{CEELA, EECA 2 y 13; indicador propio} \\[4pt]
E02 · Eficiencia energética operacional & Ahorro $\geq$ 20\% en la categoría energética EDGE. & Comparar energía operacional del diseño con el caso de referencia calculado en EDGE. El CB de la tesis no sustituye automáticamente la referencia EDGE.\par\textit{EDGE V3, Part 1, p. 7; referencia externa} \\[4pt]
E03 · Transmitancia de cerramientos & U $\leq$ límite A2 aplicable. Ejemplo sin atenuación, zonas 0–2: cubierta ligera 0,40; muro sólido 0,75. & U del sistema completo y límite por elemento, clima, tipología constructiva y compacidad. Confirmar zona climática y compacidad; Ug y proporción acristalada se verifican por fachada.\par\textit{Minergie LATAM 2025.1, A2, tabla 2, pp. 15–16} \\[4pt]
E04 · Protección solar eficaz & En zonas 0–3: $\geq$ 90\% con SHGC modificado < 0,20 durante riesgo de sobrecalentamiento. & Área transparente eficazmente protegida / área transparente total. Revisar orientación, sombras y vidrio conjuntamente; longitud del alero no acredita cumplimiento.\par\textit{Minergie LATAM, guía 2025.1, A5, pp. 35–38; CEELA, EECA 2} \\[4pt]
\bottomrule\end{tabular}}\fuente{Elaboración propia con las versiones documentales indicadas en esta sección. Los criterios propios se identifican expresamente.}\end{table}\clearpage
\subsection{Eficiencia y ambiente interior}
\begin{table}[H]\centering\caption{Indicadores mínimos: eficiencia y ambiente interior}
{\fontsize{9.5}{12}\selectfont\singlespacing\renewcommand{\arraystretch}{1.25}\setlength{\tabcolsep}{4pt}\begin{tabular}{>{\raggedright\arraybackslash}p{2.9cm}>{\raggedright\arraybackslash}p{5.2cm}>{\raggedright\arraybackslash}p{5.9cm}}\toprule Indicador & Criterio de referencia & Cálculo, alcance y fuente \\ \midrule
E05 · Potencia y control de iluminación & $\leq$ 8 W/m\textsuperscript{2}, LED y control sobre $\geq$ 80\% de potencia o todos los recintos regularmente ocupados. & Potencia instalada / área por recinto; verificar LED y alcance de automatización. Requisito para oficinas y espacios exteriores; conservar iluminancia de tarea y comprobar LED con ficha técnica.\par\textit{Minergie LATAM 2025.1, T3, p. 27; CEELA, EECA 10} \\[4pt]
C01 · Aceptabilidad térmica ocupada & Ruta mecánica seleccionada: PMV entre $-$0,5 y +0,5 en $\geq$ 90\% de horas y $\geq$ 90\% de espacios regularmente ocupados. & Evaluar PMV por recinto y hora; en modo natural aplicar modelo adaptativo admisible. En modo mixto separar periodos mecánicos y naturales; documentar met, clo, velocidad, humedad y radiación.\par\textit{WELL v2 Q4 2020, T01.1, opción 1, p. 161} \\[4pt]
C02 · Suministro de aire exterior & En los recintos indicados por T4: caudal $\geq$ 1,30 veces referencia local o ASHRAE 62.1 si no existe. & Separar aire exterior de infiltración; verificar caudal, filtración y recuperación. Aplicar variante oficinas y condiciones de exención; ventilación natural no prueba calidad de aire.\par\textit{Minergie LATAM 2025.1, T4, pp. 28–29; CEELA, EECA 6} \\[4pt]
C03 · Autonomía lumínica y exposición solar & sDA $\geq$ 40\% como nivel de entrada del crédito; documentar control de deslumbramiento donde ASE supere 10\%. & Calcular sDA300/50\% y ASE1000,250 con simulación anual. Crédito opcional usado como referencia; no es requisito mínimo para certificar todo el edificio.\par\textit{LEED v4.1 BD+C, guía abril 2019, EQ Daylight, opción 1} \\[4pt]
\bottomrule\end{tabular}}\fuente{Elaboración propia con las versiones documentales indicadas en esta sección. Los criterios propios se identifican expresamente.}\end{table}\clearpage
\subsection{Agua, materiales y generación}
\begin{table}[H]\centering\caption{Indicadores mínimos: agua, materiales y generación}
{\fontsize{9.5}{12}\selectfont\singlespacing\renewcommand{\arraystretch}{1.25}\setlength{\tabcolsep}{4pt}\begin{tabular}{>{\raggedright\arraybackslash}p{2.9cm}>{\raggedright\arraybackslash}p{5.2cm}>{\raggedright\arraybackslash}p{5.9cm}}\toprule Indicador & Criterio de referencia & Cálculo, alcance y fuente \\ \midrule
A01 · Uso eficiente del agua & Ahorro $\geq$ 20\% en la categoría agua EDGE. & Volumen de agua según usos, caudales, frecuencias y ocupación; contraste EDGE. Mantener igual población y frecuencia de uso; documentar todos los usos aplicables.\par\textit{EDGE V3, Part 1, p. 7; CEELA, EECA 12} \\[4pt]
M01 · Carbono incorporado & Meta de reducción $\geq$ 20\% para categoría materiales EDGE; en V3, cálculo de carbono incorporado. & Sumar cantidades × factores ambientales compatibles; contraste material EDGE. Fijar misma unidad funcional y alcance de módulos; añadir transporte insular por separado.\par\textit{EDGE V3, Part 5, introducción y anexo 1; CEELA, EECA 3} \\[4pt]
M02 · Emisiones de productos interiores & Meta documental propia: 100\% de productos de la intervención con verificación de emisiones; aplicar límites del ensayo seleccionado. & Inventario de acabados, adhesivos y pinturas con ficha de emisiones aplicable. La cobertura documental no equivale a obtener el crédito WELL X06.\par\textit{WELL v2 Q4 2020, X06; CEELA, EECA 5} \\[4pt]
E06 · Autoproducción & Ruta fotovoltaica: $\geq$ 0,010 kWp/m\textsuperscript{2} SRE; el requisito se limita a 30 kWp. & Dimensionar respecto a superficie de referencia energética y estimar producción. Ruta seleccionada para oficinas; justificar superficie, sombras y posibles restricciones locales.\par\textit{Minergie LATAM 2025.1, T2, p. 26; CEELA, EECA 14} \\[4pt]
\bottomrule\end{tabular}}\fuente{Elaboración propia con las versiones documentales indicadas en esta sección. Los criterios propios se identifican expresamente.}\end{table}\clearpage
\subsection{Sistemas, operación y adaptación insular}
\begin{table}[H]\centering\caption{Indicadores mínimos: sistemas, operación y adaptación insular}
{\fontsize{9.5}{12}\selectfont\singlespacing\renewcommand{\arraystretch}{1.25}\setlength{\tabcolsep}{4pt}\begin{tabular}{>{\raggedright\arraybackslash}p{2.9cm}>{\raggedright\arraybackslash}p{5.2cm}>{\raggedright\arraybackslash}p{5.9cm}}\toprule Indicador & Criterio de referencia & Cálculo, alcance y fuente \\ \midrule
E07 · Refrigeración eficiente & SEER $\geq$ 6 kW/kW (método europeo) y sobredimensionamiento $\leq$ 15\% en oficinas; revisar refrigerante y recuperación. & Verificar ficha de equipo, refrigerante y potencia instalada frente a carga de diseño. No convertir directamente COP del modelo en SEER; aplicar las condiciones climáticas de T5.\par\textit{Minergie LATAM 2025.1, T5, p. 30; CEELA, EECA 13} \\[4pt]
G01 · Monitoreo y control & Plan completo de energía, agua y ambiente interior para la fase de operación. & Comprobar plan con variables, medidores, responsables, frecuencia y mantenimiento. Evaluación de diseño; los datos simulados no sustituyen medición posterior.\par\textit{CEELA, EECA 15; Minergie O1–O3; criterio documental propio} \\[4pt]
G02 · Viabilidad de materiales y soluciones & Todas las soluciones propuestas con ficha de viabilidad y fuentes; menor impacto debe demostrarse. & Comparar procedencia, transporte, vida útil, humedad/salinidad, mantenimiento y costo. Criterio propio de selección; no supone que cualquier material local o reciclado sea automáticamente sostenible.\par\textit{CEELA, EECA 1 y 3; adaptación propia a Galápagos} \\[4pt]
P01 · Apertura útil y enfriamiento natural & Si se emplea enfriamiento nocturno natural: aberturas $\geq$ 4\% del piso en oficinas, protegidas de lluvia y robos. & Abertura neta disponible / área de piso del recinto; comprobar recorrido y horario. Criterio condicionado a esa estrategia; fracción practicable de ventana y proporción respecto al piso son diferentes.\par\textit{Minergie LATAM 2025.1, A6, pp. 18–19; CEELA, EECA 6 y 8} \\[4pt]
\bottomrule\end{tabular}}\fuente{Elaboración propia con las versiones documentales indicadas en esta sección. Los criterios propios se identifican expresamente.}\end{table}\clearpage
\subsection{Aplicación y validación comparativa}
La matriz se aplica primero al caso base (sección~\ref{sec:matriz-caso-base}) y después a las alternativas (sección~\ref{sec:matriz-escenarios}). Para cada corrida se conservan clima, calendario, ocupación, cargas de referencia y alcance de zonas. Las variables de intervención se registran antes de interpretar resultados. Si una revisión modifica un supuesto del CB, deben recalcularse las alternativas que se comparen con esa versión.

Para una magnitud cuyo descenso es favorable, la mejora porcentual se define como $100(X_{CB}-X_i)/X_{CB}$. En confort se reporta también la diferencia en puntos porcentuales de horas aceptables: $H_i-H_{CB}$. Si el denominador es cero o faltan datos, no se informa un porcentaje. Los porcentajes mensuales no se promedian para obtener el anual.

La evaluación térmica se realiza por recinto regularmente ocupado y horario efectivo. En espacios de modo mixto se aplica el método correspondiente al modo activo en cada periodo. Se documentan vestimenta, actividad, humedad, temperatura radiante y velocidad del aire. Las medias del edificio y el recuento de días sobre 28 °C se conservan solo como diagnóstico exploratorio. No se equipara una reducción de enfriamiento a mejora del confort.

La comparación de materiales exige igual función, vida útil y límites del cálculo ambiental. Una mejora de reflectancia o transmitancia puede reducir demanda térmica, pero no demuestra menor carbono incorporado. La validación de diseño comprende consistencia del modelo, trazabilidad de datos y sensibilidad de supuestos; sin mediciones de operación no se declara calibración empírica. La revisión documental de la matriz queda abierta a contraste técnico externo antes de proponer su uso generalizado.
